\chapter{Ejemplos OpenAPI y BIAN}
\label{anexo:f}

Este anexo complementa la Tabla~1.4 (Cap.~1 §1.4) con un \textbf{ejemplo reducido} y un
\textbf{ejemplo ampliado} de contrato OpenAPI bancario contrastado con un \ServiceDomain{} BIAN
genérico del dominio de pagos. Los ejemplos son \textbf{ilustrativos}: no fijan un caso piloto
único.

\section{Ejemplo reducido (unidad de análisis)}

\begin{table}[H]
  \centering
  \caption{Ejemplo reducido — pares OpenAPI $\leftrightarrow$ BIAN}
  \label{tab:anexo-f-reducido}
  \small
  \begin{tabular}{@{}p{0.44\textwidth}p{0.44\textwidth}@{}}
    \toprule
    \textbf{Elemento OpenAPI} & \textbf{Elemento BIAN (referencia)} \\
    \midrule
    \texttt{POST /payment-orders} & \Behavior{} \texttt{Initiate} \\
    \texttt{GET /payment-orders/\{id\}} & \Behavior{} \texttt{Retrieve} \\
    Esquema \texttt{PaymentOrder} (\texttt{amount}, \texttt{currency}) &
    \BusinessObject{} \texttt{Payment} / \texttt{PaymentAmount} \\
    \bottomrule
  \end{tabular}
\end{table}

\section{Ejemplo ampliado (fragmento de contrato)}

\subsection{Contrato OpenAPI (extracto ilustrativo)}

\lstinputlisting{../datos-objeto-estudio/ejemplo_openapi_payment_initiation.yaml}

\subsection{Referencia BIAN (Payment Initiation — ilustrativo)}

\begin{table}[H]
  \centering
  \caption{Elementos BIAN emparejables con el extracto OpenAPI}
  \label{tab:anexo-f-bian}
  \small
  \begin{tabular}{@{}p{0.28\textwidth}p{0.62\textwidth}@{}}
    \toprule
    \textbf{Tipo BIAN} & \textbf{Elemento de referencia} \\
    \midrule
    \ServiceDomain{} & Payment Initiation (dominio de pagos) \\
    \Behavior{} & \texttt{Initiate}, \texttt{Retrieve}, \texttt{Update} \\
    \BusinessObject{} & \texttt{PaymentInstruction}, \texttt{PaymentAmount}, \texttt{Party} \\
    Relación de control & Orquestación iniciación $\rightarrow$ confirmación de estado \\
    \bottomrule
  \end{tabular}
\end{table}

\subsection{Correspondencias para el modelo S/E/C}

\begin{table}[H]
  \centering
  \caption{Lectura del ejemplo bajo S, E y C}
  \label{tab:anexo-f-scores}
  \small
  \begin{tabular}{@{}p{0.12\textwidth}p{0.78\textwidth}@{}}
    \toprule
    \textbf{Score} & \textbf{Qué evaluaría en este extracto} \\
    \midrule
    E &
    Correspondencia \texttt{/payment-initiations} $\leftrightarrow$ \Behavior{};
    anidamiento \texttt{components/schemas} $\leftrightarrow$ \BusinessObject{} \\
    S &
    Similitud \texttt{instructedAmount/currency} $\leftrightarrow$ \texttt{PaymentAmount};
    \texttt{creditor} $\leftrightarrow$ \texttt{Party} \\
    C &
    Proporción de \Behavior{} y \BusinessObject{} del SD cubiertos por operaciones y esquemas
    del fragmento \\
    \bottomrule
  \end{tabular}
\end{table}

\section{Escenarios numéricos ilustrativos}

Los valores siguientes son \textbf{ilustrativos} ($\alpha=\beta=\gamma=\frac{1}{3}$) para contrastar
tres niveles de alineación. La calibración empírica queda para Tesis~2.

\begin{table}[H]
  \centering
  \caption{Escenarios ilustrativos --- S, E, C $\rightarrow$ \AlignmentScore{}}
  \label{tab:anexo-f-escenarios}
  \small
  \begin{tabular}{@{}lccccp{0.38\textwidth}@{}}
    \toprule
    \textbf{Caso} & \textbf{E} & \textbf{S} & \textbf{C} & \textbf{Score} & \textbf{Clasificación} \\
    \midrule
    A (alta) & 0{,}88 & 0{,}85 & 0{,}90 & 0{,}88 & Alta --- Payment Initiation alineado \\
    B (baja) & 0{,}55 & 0{,}60 & 0{,}50 & 0{,}55 & Baja --- cobertura parcial del SD \\
    C (nula) & 0{,}12 & 0{,}15 & 0{,}10 & 0{,}12 & Nula --- dominios no relacionados \\
    \bottomrule
  \end{tabular}
\end{table}

\textbf{Ejemplo de E:} 7 nodos estructurales coincidentes sobre 8 esperados $\Rightarrow E = 7/8 =
0{,}875$. Ese valor mide similitud entre contratos normalizados; validar el \textit{prototipo}
requiere contrastar correspondencias con \textit{ground truth} de especialistas (§1.6).

\section{Nota metodológica}

El procedimiento completo de normalización, matching y cálculo de scores se documenta en el
Capítulo~4 (diseño metodológico). Este anexo complementa la Tabla~1.4 con un ejemplo reducido y
otro ampliado, sin comprometer la decisión de no fijar un piloto en la matriz de consistencia.
